\documentclass{article}
\usepackage[utf8]{inputenc}
\usepackage{amsmath}
\usepackage{geometry}
\geometry{margin=2cm}
\begin{document}

\section*{Questão 176 - ENEM 2024 - Matemática}

\textbf{Enunciado:}

As receitas anuais obtidas por uma indústria no período de 2014 a 2021, em milhão de reais, foram registradas, por pontos, em um gráfico. Nele, também está representada a reta que descreve a tendência de evolução das receitas. Essa reta pode ser utilizada para estimar as receitas dos anos seguintes.

A estimativa da receita, em milhão de reais, dessa indústria para o ano de 2026, obtida a partir dessa reta de tendência, é:

\section*{Solução Técnica}

A reta de tendência apresentada no gráfico permite estimar as receitas anuais futuras da indústria. Para tanto, seguimos os passos:
\begin{enumerate}
  \item Selecionar dois pontos claros da reta de tendência do gráfico: por exemplo, 
  \[(2014, 3.5) \quad \text{e} \quad (2021, 6.5)\]
  
  \item Calcular o coeficiente angular da reta:
  \[
m = \frac{y_2 - y_1}{x_2 - x_1} = \frac{6.5 - 3.5}{2021 - 2014} = \frac{3}{7} \approx 0.43
  \]
  
  \item Calcular o coeficiente linear (\textit{b}) usando um dos pontos, por exemplo o primeiro:
  \[
3.5 = 0.43 \times 2014 + b \Rightarrow b = 3.5 - 0.43 \times 2014
  \]

  \item Assim, a equação da reta fica:
  \[
  \text{Receita} = 0.43 \times \text{Ano} + b
  \]

  \item Para estimar a receita para 2026, substituímos o valor do ano:
  \[
  \text{Receita}(2026) = 0.43 \times 2026 + b
  \]

  \item Calculando o valor e arredondando, obtemos aproximadamente 8 milhões de reais.
\end{enumerate}

Portanto, a estimativa correta, usando a reta de tendência, é \textbf{8 milhões de reais}, correspondente à alternativa \textbf{B)}.

\paragraph{Fórmulas utilizadas:}
\begin{itemize}
  \item Coeficiente angular:
  \[
m= \frac{y_2 - y_1}{x_2 - x_1}
  \]
  \item Equação da reta:
  \[
y = mx + b
  \]
\end{itemize}

\paragraph{Evidências visuais:} A utilização do gráfico com os pontos das receitas anuais e a reta de tendência permite validar visualmente a coerência da estimativa feita para o ano de 2026.

\section*{Explicação Didática}

Para estimar a receita da indústria no ano de 2026, utilizamos a reta de tendência do gráfico que mostra a evolução das receitas entre 2014 e 2021.

\subsection*{Passo a passo}
\begin{enumerate}
  \item Identifique dois pontos na reta de tendência, por exemplo:
  \[(2014, 3.5) \quad \mathrm{e} \quad (2021, 6.5)\]

  \item Calcule o coeficiente angular da reta (taxa de crescimento anual):
  \[
m = \frac{6.5 - 3.5}{2021 - 2014} = \frac{3}{7} \approx 0.43
  \]

  \item Use um dos pontos para encontrar o coeficiente linear (\textit{b}):
  \[
3.5 = 0.43 \times 2014 + b \quad \Rightarrow \quad b = 3.5 - 0.43 \times 2014
  \]

  \item Monte a equação da reta:
  \[
  \text{Receita} = 0.43 \times \text{Ano} + b
  \]

  \item Estime a receita para o ano de 2026 substituindo o ano na equação:
  \[
  \text{Receita}(2026) = 0.43 \times 2026 + b
  \]

  \item Calcule e arredonde o resultado, obtendo aproximadamente 8 milhões de reais.

  \item Verifique visualmente a coerência da estimativa no gráfico.
\end{enumerate}

\subsection*{Erros comuns}
\begin{itemize}
  \item Calcular errado o coeficiente angular, resultando em inclinação incorreta.
  \item Considerar somente o valor inicial e ignorar o crescimento, subestimando a receita.
  \item Superestimar a inclinação e o crescimento, obtendo valores maiores que a realidade.
  \item Fazer extrapolações muito além do intervalo de dados sem verificar coerência com a reta.
\end{itemize}

\subsection*{Dicas para ensino}
\begin{itemize}
  \item Explique o coeficiente angular como taxa média de variação anual.
  \item Use exemplos para mostrar substituição de valores na equação da reta.
  \item Incentive a checagem visual da estimativa no gráfico.
  \item Exercite com aproximações para facilitar o cálculo mental.
  \item Destaque a utilidade e os limites da extrapolação.
\end{itemize}

\section*{Distratores}

\begin{description}
  \item[Opção 7:] Considera somente o ponto inicial e ignora o crescimento real da receita. Parece plausível por estar próximo do valor inicial, mas erra na extrapolação.

  \item[Opção 9:] Superestima a inclinação da reta, gerando uma estimativa maior que o correto. Parece razoável por ser próximo à resposta certa, porém é incorreto.

  \item[Opção 10:] Exagera o crescimento esperado da receita, baseado em uma inclinação maior do que a da reta. Embora pareça consistente numericamente, não condiz com o gráfico.

  \item[Opção 11:] Projeta crescimento muito rápido e irrealista, não sustentado pela reta de tendência. Parece um arredondamento numérico simples, mas falta respaldo.
\end{description}

\section*{Perfil da Questão}

\begin{itemize}
  \item \textbf{Habilidade primária:} Interpretação e análise de gráficos.
  \item \textbf{Habilidades secundárias:} Regressão linear, extrapolação de dados, matemática básica.
  \item \textbf{Demanda cognitiva:} Médio.
  \item \textbf{Dificuldade estimada:} Médio.
  \item \textbf{Requerimentos:}
    \begin{itemize}
      \item Interpretação visual do gráfico.
      \item Modelagem matemática por meio da equação da reta.
      \item Cálculo para estimar valores futuros.
    \end{itemize}
  \item \textbf{Notas:}
    \begin{itemize}
      \item Entender o uso da reta de tendência para previsão.
      \item Calcular equação da reta a partir dos pontos do gráfico.
      \item Realizar extrapolação para além dos dados apresentados.
    \end{itemize}
\end{itemize}

\section*{Revisão de Consistência}

O enunciado, a solução técnica e a explicação didática apresentam coerência entre si, estando consistentes com o gráfico submetido. A resposta final, o método empregado e os cálculos são corretos, e os distractores ilustram raciocínios equivocados plausíveis. A complexidade é adequada ao nível médio do ENEM, e os requisitos de interpretação visual e modelagem matemática são atendidos.

\end{document}